% ============================================================
% Chuong 3 -- Ket qua thuc tap
% Sinh vien: Dinh Van Tan Luong -- MSV: 1771020446
% Don vi: Cong ty Co Phan Tuong Lai NEXTX
% ============================================================

\chapter{KẾT QUẢ THỰC TẬP}

\section{Tóm tắt nội dung đã thực hiện}

Trong 9 tuần thực tập tại Công ty Cổ phần Tương lai NEXTX, sinh viên đã tham gia xây
dựng toàn bộ phần xử lý phía sau (backend) của hệ thống \textbf{NextFarm Chatbot Nông
Nghiệp AI (phiên bản 2.2)}, đi theo đúng 4 giai đoạn đã trình bày chi tiết ở Chương 2:
(1) nghiên cứu và xác định yêu cầu, nguồn dữ liệu, ràng buộc tích hợp NextFarm; (2) xây
dựng nền tảng kỹ thuật gồm kiến trúc Tri-Layer RAG và lớp Tool kết nối IoT; (3) xây dựng
cơ chế định tuyến (Router hai cấp) và tổng hợp câu trả lời (Output, Guardrail); (4) hoàn
thiện Knowledge Graph, tích hợp toàn bộ pipeline và kiểm thử end-to-end.

Song song với công việc chuyên môn, sinh viên thường xuyên báo cáo tiến độ hằng tuần,
trao đổi trực tiếp với người hướng dẫn tại doanh nghiệp để làm rõ yêu cầu nghiệp vụ nông
nghiệp, đồng thời phối hợp với đội kỹ thuật phụ trách nền tảng NextFarm để xác minh lại
đặc tả API trong các trường hợp tài liệu chưa mô tả rõ (đặc biệt ở giai đoạn thiết kế
ràng buộc phân quyền tại Tuần 3). Toàn bộ quá trình làm việc, trao đổi và các phiên bản
mã nguồn đều được lưu vết qua Git/GitHub, giúp dễ dàng đối chiếu tiến độ thực tế so với
kế hoạch đề ra ban đầu.

Nhìn chung, sinh viên đã hoàn thành đầy đủ khối lượng công việc được giao: từ khảo sát
dữ liệu, thiết kế kiến trúc, xây dựng mã nguồn, cho đến kiểm thử và đánh giá hệ thống
bằng bộ benchmark nội bộ, đáp ứng đúng mục tiêu đã đề ra ở Chương 2.

\section{Sản phẩm hoặc dự án thực tế đã tham gia}

\subsection{Thông tin dự án}

\begin{itemize}
    \item \textbf{Tên dự án}: NextFarm Chatbot Nông Nghiệp AI (phiên bản 2.2).
    \item \textbf{Mục tiêu dự án}: xây dựng chatbot tư vấn nông nghiệp đáng tin cậy,
    tích hợp trực tiếp với nền tảng IoT NextFarm để trả lời cả câu hỏi kiến thức canh
    tác tĩnh lẫn câu hỏi vận hành thời gian thực theo từng trang trại.
    \item \textbf{Vai trò của sinh viên}: thực tập sinh phát triển backend
    (\textit{Intern Backend Developer}), phụ trách toàn bộ chuỗi xử lý câu hỏi phía sau
    của chatbot — từ tiền xử lý ngôn ngữ, định tuyến, truy xuất đa nguồn, tổng hợp câu
    trả lời, đến cơ chế kiểm soát chất lượng (Guardrail).
\end{itemize}

\subsection{Các chức năng, nhiệm vụ đã thực hiện}

\begin{itemize}
    \item Khảo sát, phân loại 4 nhóm nguồn dữ liệu và thiết kế kiến trúc Tri-Layer RAG
    (Fact Store -- Knowledge Graph -- Document Store).
    \item Thiết kế đặc tả ràng buộc an toàn khi tích hợp API NextFarm (phân quyền theo
    trang trại, chính sách fail-closed, nguyên tắc chỉ đọc dữ liệu thiết bị).
    \item Xây dựng 6 hàm tool kết nối API NextFarm (lấy cảm biến mới nhất, trạng thái
    thiết bị, lịch tưới, lịch sử tưới, lịch sử lệnh điều khiển, cảnh báo) và các route
    REST tương ứng trong \texttt{tool\_router.py}.
    \item Xây dựng Router hai cấp: Fast-Path Rule Router (rule-based) và LLM Router
    (dùng Gemini) làm phương án dự phòng.
    \item Xây dựng lớp Output tổng hợp câu trả lời đa nguồn và module \texttt{validator.py}
    (Guardrail) chống bịa đặt dữ liệu.
    \item Xây dựng hoàn chỉnh tầng Knowledge Graph (bảng \texttt{kg\_triples}) cho 10
    loại cây trồng chính, hỗ trợ truy vấn quan hệ nhiều bước.
    \item Phát hiện và khắc phục lỗi tích hợp nghiêm trọng ở Retrieval Plan (thiếu 4/6
    hàm tool) thông qua kiểm thử end-to-end toàn bộ pipeline.
    \item Xây dựng và chạy bộ benchmark nội bộ hơn 260 câu hỏi để đánh giá độ chính xác
    và độ an toàn của hệ thống.
\end{itemize}

Bảng~\ref{tab:cong-nghe-su-dung} tổng hợp các công nghệ, công cụ chính đã sử dụng trong
suốt quá trình thực tập.

\begin{table}[H]
\centering
\caption{Danh sách công nghệ, công cụ sử dụng trong quá trình thực tập}
\label{tab:cong-nghe-su-dung}
\begin{tabular}{|c|p{4cm}|p{8cm}|}
\hline
\textbf{STT} & \textbf{Công nghệ/Công cụ} & \textbf{Mục đích sử dụng} \\
\hline
1 & Python, FastAPI & Xây dựng toàn bộ backend, API xác thực, chat và quản trị. \\
\hline
2 & Google Gemini API & Mô hình ngôn ngữ lớn dùng cho LLM Router, tổng hợp câu trả lời
và Guardrail (LLM-as-Judge). \\
\hline
3 & PostgreSQL & Lưu trữ dữ liệu có cấu trúc: bảng \texttt{facts} (số liệu định lượng)
và \texttt{kg\_triples} (Knowledge Graph). \\
\hline
4 & ChromaDB & Cơ sở dữ liệu vector lưu embedding tài liệu kỹ thuật (Document Store). \\
\hline
5 & BM25 Okapi (rank\_bm25) & Truy xuất thưa (sparse retrieval), kết hợp Dense retrieval
qua thuật toán RRF. \\
\hline
6 & underthesea & Tiền xử lý, tách từ và chuẩn hóa câu hỏi tiếng Việt. \\
\hline
7 & marker & Phân tích, trích xuất nội dung tài liệu PDF/DOCX/PPTX đầu vào. \\
\hline
8 & Git/GitHub & Quản lý mã nguồn, lưu vết tiến độ và phối hợp với người hướng dẫn. \\
\hline
9 & Visual Studio Code & Soạn thảo và gỡ lỗi mã nguồn. \\
\hline
\end{tabular}
\Nguon{Sinh viên tổng hợp}
\end{table}

\section{Kỹ năng đã rèn luyện được}

\subsection{Kỹ năng chuyên môn}

\begin{itemize}
    \item Thiết kế cơ sở dữ liệu quan hệ cho dữ liệu định lượng và đồ thị tri thức
    (schema bảng \texttt{facts}, \texttt{kg\_triples}), viết truy vấn SQL nối bảng phục
    vụ suy luận quan hệ nhiều bước.
    \item Xây dựng API REST với FastAPI, thiết kế middleware kiểm tra xác thực và phân
    quyền trước khi xử lý nghiệp vụ.
    \item Kỹ thuật truy xuất thông tin nâng cao: kết hợp Dense retrieval (embedding
    đa ngôn ngữ) và Sparse retrieval (BM25), hợp nhất kết quả bằng thuật toán RRF.
    \item Kỹ năng viết prompt (\textit{prompt engineering}) cho mô hình ngôn ngữ lớn để
    vừa trích xuất tham số có cấu trúc (JSON), vừa sinh câu trả lời tự nhiên đúng phạm
    vi dữ liệu cho phép.
    \item Kỹ năng gỡ lỗi (debug) một hệ thống pipeline nhiều tầng, biết cách kiểm thử
    end-to-end để phát hiện lỗi tích hợp mà kiểm thử từng module riêng lẻ không phát
    hiện được.
    \item Viết unit test cho các hàm xử lý nghiệp vụ, đặc biệt với các trường hợp biên
    (dữ liệu thiếu, thiết bị offline, phân quyền không hợp lệ).
    \item Kỹ năng đọc hiểu và làm việc với tài liệu API của một hệ thống bên ngoài
    (NextFarm) để thiết kế lớp tích hợp an toàn.
\end{itemize}

\subsection{Kỹ năng mềm}

\begin{itemize}
    \item Kỹ năng giao tiếp và đặt câu hỏi đúng trọng tâm khi làm việc với người hướng
    dẫn và đội kỹ thuật NextFarm để làm rõ các điểm chưa rõ trong tài liệu API.
    \item Kỹ năng quản lý thời gian và lập kế hoạch công việc theo từng tuần, bám sát
    tiến độ 4 giai đoạn của dự án trong suốt 9 tuần.
    \item Kỹ năng viết tài liệu kỹ thuật rõ ràng, mạch lạc (đặc tả ràng buộc tích hợp,
    nhật ký thực tập), giúp người khác dễ dàng nắm bắt lại quyết định thiết kế.
    \item Tư duy phản biện và cẩn trọng khi rà soát hệ thống, thể hiện rõ qua việc chủ
    động kiểm thử end-to-end và phát hiện lỗi tích hợp nghiêm trọng ở Tuần 9 thay vì chỉ
    dừng lại khi từng module riêng lẻ đã chạy đúng.
    \item Ý thức trách nhiệm với chất lượng sản phẩm khi hiểu rằng người dùng cuối của
    hệ thống là nông dân thực tế — một câu trả lời sai về liều lượng phân bón có thể gây
    thiệt hại kinh tế trực tiếp, từ đó hình thành thái độ làm việc cẩn trọng, luôn ưu
    tiên sự an toàn và trung thực của câu trả lời hơn là tính ``mượt mà''.
\end{itemize}

\section{Đề xuất cải thiện hoặc đóng góp ý kiến cho doanh nghiệp}

Dựa trên những khó khăn và bài học thực tế trong quá trình thực tập (đã trình bày chi
tiết ở Chương 2), sinh viên xin đề xuất một số ý kiến sau:

\begin{itemize}
    \item \textbf{Bổ sung bước kiểm thử end-to-end bắt buộc}: nên đưa kiểm thử xuyên
    suốt toàn bộ pipeline thành bước bắt buộc trước khi nghiệm thu bất kỳ module nào,
    thay vì chỉ kiểm thử đơn vị (unit test) riêng lẻ — tránh lặp lại tình huống như ở
    Tuần 9, khi lỗi Retrieval Plan thiếu tool chỉ được phát hiện ở giai đoạn tích hợp
    cuối cùng.
    \item \textbf{Đảm bảo benchmark phản ánh đúng pipeline thực tế}: công cụ đánh giá
    nên luôn gọi qua đúng luồng xử lý sản xuất (lần lượt qua Router, Retrieval Plan rồi
    đến bước tổng hợp câu trả lời) thay vì gọi trực tiếp vào từng hàm xử lý, để tránh
    đánh giá ``ảo'' không phản ánh đúng chất lượng thực tế của hệ thống.
    \item \textbf{Hoàn thiện thêm tài liệu API nội bộ}: một số điểm trong tài liệu API
    của NextFarm (đặc biệt về cơ chế phân quyền) chưa được mô tả đầy đủ, khiến thực tập
    sinh mất thời gian trao đổi lại với đội kỹ thuật; nên cập nhật tài liệu API đầy đủ
    và dễ tiếp cận hơn cho các thực tập sinh ở các đợt sau.
    \item \textbf{Tổ chức buổi rà soát kiến trúc giữa kỳ}: nên bố trí một buổi review
    kiến trúc tổng thể ở khoảng giữa thời gian thực tập (thay vì chỉ ở giai đoạn tích
    hợp cuối) để phát hiện sớm hơn các điểm chưa khớp giữa các module do nhiều người/
    giai đoạn phát triển riêng lẻ.
    \item \textbf{Tiếp tục đầu tư cho module Guardrail}: do đây là lớp quyết định độ an
    toàn của toàn hệ thống khi phục vụ người dùng thật, doanh nghiệp nên dành thêm thời
    gian và dữ liệu kiểm thử cho nhóm câu hỏi biên (ngoài phạm vi, không có đáp án, cố
    tình truy cập trái phép) ở các giai đoạn phát triển tiếp theo.
\end{itemize}

\section{Tự đánh giá hiệu quả thực tập}

Qua 9 tuần thực tập, sinh viên tự đánh giá đã hoàn thành đầy đủ mục tiêu và khối lượng
công việc được giao: xây dựng thành công toàn bộ chuỗi xử lý backend của chatbot nông
nghiệp, từ khâu khảo sát dữ liệu đến kiểm thử, đánh giá hệ thống bằng bộ benchmark nội
bộ.

\textbf{Điểm mạnh}: sinh viên nắm vững và vận dụng được một kiến trúc truy xuất thông
tin nâng cao (Tri-Layer RAG kết hợp Agentic Tools), có khả năng tự nghiên cứu tài liệu
kỹ thuật mới (API bên thứ ba, kỹ thuật hybrid retrieval) và áp dụng vào bài toán thực
tế; đồng thời thể hiện được tinh thần trách nhiệm khi chủ động kiểm thử kỹ lưỡng và phát
hiện lỗi hệ thống trước khi bàn giao.

\textbf{Hạn chế}: ở giai đoạn đầu, do chưa có nhiều kinh nghiệm về nghiệp vụ nông
nghiệp, sinh viên mất khá nhiều thời gian tìm hiểu thuật ngữ và quy trình canh tác trước
khi có thể thiết kế dữ liệu chính xác; bên cạnh đó, việc chỉ kiểm thử từng module riêng
lẻ ở các tuần đầu khiến lỗi tích hợp ở lớp Retrieval Plan không được phát hiện sớm, dẫn
đến phải dành thêm thời gian khắc phục ở giai đoạn cuối.

\textbf{Định hướng sau thực tập}: sinh viên mong muốn tiếp tục trau dồi kiến thức và kỹ
năng thiết kế hệ thống AI theo hướng an toàn, đáng tin cậy (\textit{safety-first AI
design}) — đặc biệt là các kỹ thuật chống bịa đặt thông tin (hallucination) cho mô hình
ngôn ngữ lớn, vì đây là năng lực ngày càng quan trọng khi triển khai AI vào các lĩnh vực
có ảnh hưởng trực tiếp đến đời sống như nông nghiệp. Nhìn chung, đợt thực tập đã giúp
sinh viên hiểu rõ hơn về quy trình phát triển một sản phẩm AI ứng dụng trong môi trường
doanh nghiệp thực tế, từ khâu thiết kế, xây dựng, kiểm thử cho đến đánh giá chất lượng
trước khi đưa vào sử dụng.